%=========== ΛΥΣΗ - 66ο ΘΕΜΑ Δ ===========
\begin{thema}{Δ}
\begin{erwthma}
\item Η $f$ είναι γνησίως αύξουσα στο $[0,5]$. Άρα είναι $1-1$ στο
$[0,5]$, οπότε αντιστρέφεται.

Επειδή η $f$ είναι συνεχής ως παραγωγίσιμη και γνησίως αύξουσα, έχουμε
\[
f([0,5])=[f(0),f(5)]=[0,4].
\]
Άρα το πεδίο ορισμού της αντίστροφης είναι
\[
{D_{f^{-1}}=[0,4]}.
\]

\item Το σκιασμένο εμβαδόν είναι
\[
E=\int_0^5 f(x)\d x.
\]
Από τα δεδομένα είναι
\[
\int_0^5 f(x)\d x=8.
\]

Η $f$ είναι γνησίως αύξουσα από το σημείο $(0,0)$ μέχρι το σημείο
$(5,4)$. Αν στο ορθογώνιο
\[
[0,5]\times[0,4]
\]
σχεδιάσουμε τη $C_f$ και τη συμμετρική της ως προς την ευθεία $y=x$,
δηλαδή την $C_{f^{-1}}$, τότε τα δύο αντίστοιχα εμβαδά συμπληρώνουν
το εμβαδόν του ορθογωνίου.

Το ορθογώνιο έχει εμβαδόν
\[
5\cdot4=20.
\]
Επομένως
\[
\int_0^5 f(x)\d x+\int_0^4 f^{-1}(x)\d x=20.
\]
Άρα
\[
8+\int_0^4 f^{-1}(x)\d x=20,
\]
οπότε
\[
{\int_0^4 f^{-1}(x)\d x=12}.
\]

\item Η $f$ είναι συνεχής στο $[0,5]$ και παραγωγίσιμη στο $(0,5)$.
Άρα, από το Θεώρημα Μέσης Τιμής, υπάρχει $\xi\in(0,5)$ τέτοιο ώστε
\[
f'(\xi)=\frac{f(5)-f(0)}{5-0}
=\frac{4-0}{5}
=\frac45.
\]
Επομένως
\[
{\exists\,\xi\in(0,5):\ f'(\xi)=\frac45}.
\]

Έστω τώρα ότι η $f$ είναι κυρτή στο $[0,5]$. Η εφαπτομένη της $C_f$
στο σημείο $(5,4)$ έχει εξίσωση
\[
\varepsilon:\quad y-4=f'(5)(x-5).
\]
Αν τέμνει τον άξονα $x'x$ στο σημείο $A(a,0)$, τότε
\[
0-4=f'(5)(a-5).
\]
Άρα
\[
a=5-\frac{4}{f'(5)}.
\]

Θα δείξουμε ότι $a>0$, δηλαδή ότι
\[
f'(5)>\frac45.
\]
Επειδή η $f$ είναι κυρτή, η γραφική της παράσταση βρίσκεται κάτω από
τη χορδή που ενώνει τα σημεία $(0,0)$ και $(5,4)$. Η χορδή αυτή έχει
εξίσωση
\[
y=\frac45x.
\]
Αν είχαμε $f'(5)=\dfrac45$, τότε η εφαπτομένη στο $(5,4)$ θα ήταν
\[
y-4=\frac45(x-5)
\Leftrightarrow y=\frac45x,
\]
δηλαδή θα συνέπιπτε με τη χορδή.

Για κυρτή συνάρτηση όμως η $C_f$ βρίσκεται πάνω από κάθε εφαπτομένη
και κάτω από κάθε χορδή. Άρα, αν η εφαπτομένη στο $(5,4)$ συνέπιπτε
με τη χορδή, θα είχαμε
\[
f(x)=\frac45x,\quad x\in[0,5].
\]
Τότε
\[
\int_0^5 f(x)\d x=\int_0^5\frac45x\d x=10,
\]
που αντιφάσκει με το δεδομένο
\[
\int_0^5 f(x)\d x=8.
\]
Άρα
\[
f'(5)>\frac45.
\]
Συνεπώς
\[
a=5-\frac{4}{f'(5)}>5-\frac{4}{\frac45}=0.
\]
Επομένως η εφαπτομένη της $C_f$ στο σημείο $(5,4)$ τέμνει τον άξονα
$x'x$ σε σημείο με θετική τετμημένη.

\item Η χορδή που ενώνει τα άκρα της $C_f$ είναι η ευθεία που διέρχεται
από τα σημεία
\[
A(0,0)\quad\text{και}\quad B(5,4).
\]
Ο συντελεστής διεύθυνσής της είναι
\[
\lambda=\frac{4-0}{5-0}=\frac45,
\]
οπότε η εξίσωσή της είναι
\[
{y=\frac45x}.
\]

Αφού η $f$ είναι κυρτή, η γραφική της παράσταση βρίσκεται κάτω από τη
χορδή $AB$. Άρα, για κάθε $x\in(0,5)$,
\[
f(x)\leq \frac45x.
\]
Επιπλέον, η ισότητα δεν μπορεί να ισχύει σε όλο το $[0,5]$, γιατί τότε
το εμβαδόν κάτω από την $C_f$ θα ήταν ίσο με $10$, ενώ είναι $8$.
Επομένως
\begin{align*}
\int_0^5 f(x)\d x
&<\int_0^5 \frac45x\d x\\
&=\left.\frac{2}{5}x^2\right|_0^5\\
&=10.
\end{align*}
Δηλαδή
\[
{\int_0^5 f(x)\d x<10}.
\]
Αυτό συμφωνεί με το δεδομένο της άσκησης, αφού
\[
\int_0^5 f(x)\d x=E=8<10.
\]
\end{erwthma}
\end{thema}
%=========== ΤΕΛΟΣ ΛΥΣΗΣ - 66ο ΘΕΜΑ Δ ===========
